\section{Read, search, and write sequences}
\label{sec:appendix-quick-start}

This example stores DNA beside ordinary metadata, reads a small part of each
sequence, and searches the stored targets. The examples assume that the
extension's installation schema is \texttt{public}. Sequence types and helper
functions use the fixed \texttt{bioseq} schema.

\subsection{Install the extension and create a table}

Install the extension in the database where the application will run. The
catalog query confirms its version in that database.

\par\noindent\begin{minipage}{\linewidth}
\begin{lstlisting}[style=bioseq-appendix-sql]
CREATE EXTENSION bioseq;
SELECT extname, extversion
FROM pg_extension
WHERE extname = 'bioseq';

CREATE TABLE sample_reads (
    read_id bigint GENERATED ALWAYS AS IDENTITY PRIMARY KEY,
    sample_id text NOT NULL,
    collected_on date NOT NULL,
    passed_qc boolean NOT NULL DEFAULT false,
    sequence bioseq.seq('dna4', 64) NOT NULL
) USING bioseq;
\end{lstlisting}
\end{minipage}

\begin{itemize}
  \item \texttt{dna4} binds the sequence to the symbols A, C, G, and T.
    Ordinary SQL input containing N fails for this alphabet.
  \item \texttt{64} is the width of a packed storage unit, in bits. It is
    neither a length limit nor the number of bases in the sequence.
  \item \texttt{USING bioseq} selects the sequence-aware table access method.
    A heap table can hold sequence values, but it does not provide the same
    table storage and custom search route.
\end{itemize}

A build may normalize accepted storage widths to 64 bits. Inspect the actual
column declaration when comparing widths; the alphabet remains unchanged.
The \texttt{psql} command \texttt{\textbackslash d+ sample\_reads} also shows
the table's access method.

\par\noindent\begin{minipage}{\linewidth}
\begin{lstlisting}[style=bioseq-appendix-sql]
SELECT a.attname AS column_name,
       format_type(a.atttypid, a.atttypmod) AS active_type
FROM pg_attribute AS a
WHERE a.attrelid = 'sample_reads'::regclass
  AND a.attnum > 0 AND NOT a.attisdropped
ORDER BY a.attnum;
\end{lstlisting}
\end{minipage}

\subsection{Choose a write policy}

For ordinary application writes, enable insert maintenance in the writing
transaction. This switch permits maintenance of indexes whose table policy
requests updates on insertion. It does not change a manually maintained
index's policy. Section~\ref{sec:appendix-indexes} lists those policies and
their interaction with the switch.

\par\noindent\begin{minipage}{\linewidth}
\begin{lstlisting}[style=bioseq-appendix-sql]
BEGIN;
SET LOCAL bioseq.index.insert_maintenance.enabled = on;
INSERT INTO sample_reads
    (sample_id, collected_on, passed_qc, sequence)
VALUES
    ('sample-a', DATE '2026-08-01', true,
     'TCACTGTAGGATCGTAGCTAGTCCGTACTGTAGCTGATCGTAGC'
     'CCTAGTCACTGTAGCCTGATCGTAGCTAGTCACTGTAGCCATACCTGAA'),
    ('sample-a', DATE '2026-08-01', true,
     'CCCCAACTTAGTCACTTCTCAAAAATTACTCTCTCATCATCATC'
     'TACTCCTAGTCACTGTAGCCTGATCGTAGCTAGTCACACACACCCCAAA'),
    ('sample-b', DATE '2026-08-02', false,
     'GGGGGTCACTTCTCAAAAATTAGGGGGGGGGTAGTCACTGTAG'
     'CCTGATCGTAGCTAGGACACATATCATCATCTACTCCTATCATCATCTACTCCTA')
RETURNING read_id, sample_id, bioseq.len(sequence) AS bases;
COMMIT;
\end{lstlisting}
\end{minipage}

Adjacent SQL string literals separated by a newline form one value in this
listing. \texttt{SET LOCAL} lasts until the transaction ends. Use ordinary
\texttt{SET} when the same policy should apply to later transactions in the
connection.

For a bulk import, turn maintenance off while loading, then build the required
index once. The following call rebuilds the minimizers-per-fragment index for
the relation.

\par\noindent\begin{minipage}{\linewidth}
\begin{lstlisting}[style=bioseq-appendix-sql]
SET bioseq.index.insert_maintenance.enabled = off;
-- Load rows with COPY, application batches, or bioseq.load.
SELECT bioseq.update_index(
    'sample_reads'::regclass, 'minimizers_per_fragment'
);
RESET bioseq.index.insert_maintenance.enabled;
\end{lstlisting}
\end{minipage}

The other minimizer recipe is \texttt{minimizers\_per\_sequence}. Import
procedures should record the table, recipe, and result of the rebuild call.
If the application requires a ready stored index, publish the dataset only
after the rebuild succeeds. A staging table allows the load and rebuild to
finish before a short application-level publication transaction. The
experimental grouped-SoA mode rejects ordinary per-tuple minimizer
maintenance. There is a streamed-load exception: \texttt{bioseq.load} can
rebuild eligible update-on-insert indexes after flushing the complete load,
subject to the grouped-format requirements in
Section~\ref{sec:appendix-indexes}.

\subsection{Read lengths, counts, and previews}

Read sequence metadata or a slice when that is all the client needs. The
following query returns lengths and symbol counts without returning the
complete sequence as text.

\par\noindent\begin{minipage}{\linewidth}
\begin{lstlisting}[style=bioseq-appendix-sql]
SELECT read_id, sample_id,
       bioseq.len(sequence) AS symbols,
       bioseq.unit_len(sequence) AS packed_units,
       bioseq.symbol_count(sequence, 'A') AS a_bases
FROM sample_reads
WHERE passed_qc
ORDER BY read_id;

SELECT read_id,
       sequence[:12]::text AS first_twelve,
       sequence[-12:]::text AS last_twelve
FROM sample_reads
ORDER BY read_id;
\end{lstlisting}
\end{minipage}

\begin{itemize}
  \item \texttt{len} returns the number of symbols; \texttt{unit\_len}
    returns the number of packed units. Both return \texttt{bigint}.
  \item \texttt{symbol\_count} counts the requested symbol. Its ability to
    use stored counts depends on the available index and selected plan.
  \item Slice bounds are zero-based and half-open: \texttt{[:12]} selects
    positions 0 through 11. \texttt{[-12:]} selects up to twelve final
    symbols. Short values are clamped to their available range.
\end{itemize}

Convert the complete value only when it is needed, for example
\texttt{SELECT sequence::text FROM sample\_reads WHERE read\_id = 1}.
Sections~\ref{sec:appendix-sequence-type}
and~\ref{sec:appendix-sequence-transforms} give the full rules for widths,
negative bounds, and transformations.

\subsection{Search targets and return alignment rows}

Put the stored target first in \texttt{mmsearch}, followed by a query with a
compatible bound alphabet. Ordinary metadata filters can remain beside the
search predicate.

\par\noindent\begin{minipage}{\linewidth}
\begin{lstlisting}[style=bioseq-appendix-sql]
SELECT read_id, sample_id,
       bioseq.query_start(read_id) AS query_start,
       bioseq.query_end(read_id) AS query_end,
       bioseq.target_start(read_id) AS target_start,
       bioseq.target_end(read_id) AS target_end,
       bioseq.cigar(read_id) AS cigar
FROM sample_reads
WHERE sample_id = 'sample-a'
  AND collected_on >= DATE '2026-08-01'
  AND passed_qc
  AND public.mmsearch(
      sequence,
      'CTAGTCACTGTAGCCTGATCGTAGCTAGT'::bioseq.seq('dna4', 64),
      public.aligner_options(best_n => 3)
  )
ORDER BY read_id, target_start, target_end;
\end{lstlisting}
\end{minipage}

The coordinate functions read the current alignment. Their \texttt{read\_id}
argument identifies the table alias owning this one search; its numeric value
is not used to find an alignment in a second table. A stored row can therefore
appear several times with different coordinates. Coordinates are zero-based
and half-open, so an interval's length is \texttt{end - start}.

The optional object changes the call's settings. Here \texttt{best\_n}
limits retained secondary regions when secondary filtering is active. It
does not limit matching table rows or all primary regions. The two-argument
form uses session defaults. Section~\ref{sec:appendix-search-options}
explains every object field, its corresponding default, and the restrictions
on constant options.

\subsection{Check the chosen plan}

Start with plain \texttt{EXPLAIN}, which does not execute the statement.
Explain the same filters and projections used by the application.

\par\noindent\begin{minipage}{\linewidth}
\begin{lstlisting}[style=bioseq-appendix-sql]
EXPLAIN (VERBOSE, COSTS OFF)
SELECT read_id,
       bioseq.len(sequence),
       bioseq.symbol_count(sequence, 'A')
FROM sample_reads
WHERE passed_qc;

EXPLAIN (VERBOSE, COSTS OFF)
SELECT read_id
FROM sample_reads
WHERE passed_qc
  AND public.mmsearch(
      sequence,
      'CTAGTCACTGTAGCCTGATCGTAGCTAGT'::bioseq.seq('dna4', 64)
  );
\end{lstlisting}
\end{minipage}

For the search, check for \texttt{Custom Scan (BioseqScan)}. The plan's
projected columns, metadata reads, index-read requirement, operand evaluation,
and pipeline describe what this statement requests. Highway dispatch fields
describe CPU implementations available to that plan. Adding
\texttt{ANALYZE} to \texttt{EXPLAIN} runs the statement and adds execution
information; it also runs writes or remote requests when those are part of the
statement.

If the search has no owning custom scan, check the table's access method,
operand types, visible top-level predicate, and enabled custom-scan path.
Then check the index state against the selected read policy. Missing index
coverage can permit recomputation under an allowing policy; it does not by
itself mean that a custom scan is impossible. The ordinary SQL function is
not an alternative row-by-row search implementation.
Section~\ref{sec:appendix-explain} explains the plan fields in detail.

\subsection{Update, delete, and bind application parameters}

Update metadata and sequence values in the same SQL statement. For indexes
maintained on insertion, the write switch also applies to the replacement
row created by an update.

\par\noindent\begin{minipage}{\linewidth}
\begin{lstlisting}[style=bioseq-appendix-sql]
BEGIN;
SET LOCAL bioseq.index.insert_maintenance.enabled = on;
UPDATE sample_reads
SET passed_qc = true,
    sequence = 'ACGTACGTACGTACGTACGTACGTACGTACGT'
WHERE read_id = 3;
COMMIT;

DELETE FROM sample_reads WHERE read_id = 3;
\end{lstlisting}
\end{minipage}

After maintenance-off inserts or sequence updates, rebuild the required
stored index before relying on its coverage. Deletion changes row visibility;
it does not require applications to identify internal index chunks themselves.

Applications should bind query text as data. The following prepared statement
keeps the sample name and query sequence separate from SQL text. The query
parameter is bound to the target alphabet before search execution.

\par\noindent\begin{minipage}{\linewidth}
\begin{lstlisting}[style=bioseq-appendix-sql]
PREPARE find_sample_reads(text, text) AS
SELECT read_id, sample_id,
       bioseq.target_start(read_id) AS target_start,
       bioseq.target_end(read_id) AS target_end
FROM sample_reads
WHERE sample_id = $1
  AND public.mmsearch(sequence, $2::bioseq.seq('dna4', 64))
ORDER BY read_id, target_start, target_end;

EXECUTE find_sample_reads(
    'sample-a', 'CTAGTCACTGTAGCCTGATCGTAGCTAGT'
);
DEALLOCATE find_sample_reads;
\end{lstlisting}
\end{minipage}

Drop \texttt{sample\_reads} when the exercise is finished.
Keep the write policy, index recipe, and
effective search options explicit in an application, and confirm the actual
plan before measuring representative searches.
